% !TEX root = main.tex
\section{Introduction}
\label{sec:intro}

A randomized program can be run by choosing a deterministic program at
random and running it. For first-order while programs, this gives the same
semantics as linear operators on measures~\cite{kozen1981}. We ask
whether it holds at higher types.

Probabilistic coherence spaces (PCS) are a model of linear logic in which a type
denotes a set of weight vectors on a set of points, its
web~\cite{girard2004,danos-ehrhard2011}. They are fully abstract for
probabilistic PCF, the extension of PCF (Programming Computable
Functions) with probabilistic
choice~\cite{ehrhard-pagani-tasson2018}. The proof of full abstraction
uses the form of the morphisms as power series, and not the definability
of every element of a space. In a concurrent game model of probabilistic
PCF, definability holds for the strategies that call one argument at a
time and depend only on the history they observe~\cite{castellan2018}.
We ask which elements of a space are mixtures of its elements with values in
$\{0,1\}$. Full abstraction does not settle this question, since it
concerns the equivalence of programs and not the laws valid on a whole
space. Such laws bound the behaviour of every program of a type at once. They
can fail on a space that contains more than the mixtures of deterministic
programs.

We call these elements the \emph{certain states} of a type. They are its
deterministic elements, and the set of their mixtures is their convex
hull, so we call it the \emph{deterministic hull} of the type. This is the construction of
generalized probabilism~\cite{paris2001,williams2012}, applied to the
points of a web. We work in multiplicative--additive linear logic (MALL)
over finite data types, whose programs use each argument exactly once. At
every MALL formula the deterministic hull lies inside the PCS. Each certain state is a clique of
Girard's coherence space. A
\emph{gap} is an element of the PCS outside the deterministic
hull. It violates a linear inequality that every certain state
satisfies~\cite{definetti1974,paris2001}.

The PCS also has an upper bound defined from certain states.
The \emph{local relaxation} of a formula is the set of weight vectors with
total weight at most $1$ on the support of every certain state of its
negation. It contains the PCS. Its orthogonal is the deterministic hull
of the negation. So every PCS lies between two sets defined from certain
states, and Table~\ref{tab:sandwich} shows where.

\begin{table}[t]
\centering
\caption{The PCS $P(A)$ between the deterministic hull and the
local relaxation of $A$. Strictness needs data types with at least two
elements, and in the last row it is shown over $\mathrm{Bool}$.}
\label{tab:sandwich}
\begin{tabular}{@{}llll@{}}
\toprule
$A$ & hull and $P(A)$ & $P(A)$ and relaxation & Sections \\
\midrule
$\mathbf{n} \multimap \mathbf{m}$ (first order) & equal & equal & \ref{sec:linear}, \ref{sec:duality-tight} \\
$T$ (second order) & strictly smaller & equal & \ref{sec:second}--\ref{sec:other}, \ref{sec:duality-local} \\
dual of $T$ & equal & strictly smaller & \ref{sec:exact-perfect}, \ref{sec:duality-bell} \\
$(T \otimes T)^{\perp}$ & strictly smaller & strictly smaller & \ref{sec:duality-shannon} \\
\bottomrule
\end{tabular}
\end{table}

At first order the three sets coincide: a probabilistic program from one
data type to another is a mixture of deterministic
programs~\cite{levin-peres-wilmer2009}. At the type
$T$ of a program that calls two functions once each, the deterministic
hull is the set of randomized sequential programs, by Kuhn's
theorem~\cite{kuhn1953}. It is the stable set polytope of a graph whose
clique-constrained polytope is the PCS. The PCS
is the whole local relaxation. When the data types have at least two
elements, the function with value $\tfrac12$ on five points of the web lies in the PCS and violates the
Klyachko--Can--Binicio\u{g}lu--Shumovsky (KCBS)
inequality~\cite{klyachko2008}, which every randomized sequential program
satisfies. So every law valid on the PCS of $T$ holds for
sequential programs, but not conversely.

At the dual of $T$, the tensor of the two argument types, the picture is
reversed. The PCS is the deterministic hull. When the data types have
at least two elements, it lies strictly inside the local relaxation. The Popescu--Rohrlich
box~\cite{popescu-rohrlich1994} lies in the local relaxation and outside
the PCS, so this gap is Bell's theorem. The two gaps are equivalent by the duality above. The deterministic hulls of a formula
and of its negation form a tight pair, each the orthogonal of the other,
iff neither has a gap. Tightness holds at first order and can fail at a tensor of first-order
types.

For the dual of a tensor of two formulas built from data types with
$\mathbin{\&}$ and $\oplus$, the gap is governed by induced four-cycles of
the coherence graphs. The direction without a gap has an explicit
decomposition, mechanized for formulas built from the unit type without
perfect graph theory. Coordinate retracts, embeddings of webs that
preserve the PCS in both directions, carry the gap to every
formula over nonempty data types in which such a tensor occurs
negatively. In quantum event logic no state has values in
$\{0,1\}$~\cite{kochen-specker1967}, so only the local relaxation
applies.

\paragraph{Contributions.}
\begin{enumerate}
  \item \textbf{Cliques.} Every certain state of a MALL formula is a clique
  of its coherence graph, but not conversely (Theorems~\ref{thm:mall}
  and~\ref{thm:shannon}).
  \item \textbf{Second order.} At $T$ the deterministic hull and the
  PCS are the stable set and the clique-constrained polytopes
  of one graph, with an explicit gap for two functional arguments
  (Theorems~\ref{thm:pentagon} and~\ref{thm:qstab}).
  \item \textbf{Exact condition.} For formulas built from data types with
  $\mathbin{\&}$ and $\oplus$, the dual of a tensor has no gap iff one
  factor has no induced four-cycle, with a direct proof mechanized over the
  unit type (Theorem~\ref{thm:second-exact}).
  \item \textbf{Transfer.} Coordinate retracts carry the gap to every
  formula over nonempty data types with such a tensor occurring negatively,
  a condition not necessary beyond second order (Theorem~\ref{thm:c4},
  Proposition~\ref{prop:f3}).
  \item \textbf{Duality.} Every PCS lies between the deterministic hull
  and the local relaxation, and the deterministic hulls of a formula and
  of its negation are tight iff neither has a gap, which holds at first
  order and can fail at a tensor of first-order types
  (Theorems~\ref{thm:loc-linear} and~\ref{thm:duality},
  Corollary~\ref{cor:tensor}).
\end{enumerate}

Every result of the paper is mechanized in Lean~4 over
\textsf{mathlib}~\cite{moura-ullrich2021,mathlib2020}, in about
$8{,}300$ lines, except steps the text marks as not formalized. A superscript number such as
\textsuperscript{\textsf{\scriptsize L1}} links a statement to its numbered
Lean counterpart. The archive in the data availability statement lists
the Lean statements.
